\item \textbf{Aproximación de \(\sin(x)\) mediante serie de Taylor}.

\begin{tcolorbox}[colback=blue!2!white,colframe=blue!55!black,title={Enunciado}]
La función seno \(\sin(x)\) se puede aproximar mediante su serie de Taylor:
\[
\sin(x) = \sum_{k=0}^{\infty} (-1)^k \frac{x^{2k+1}}{(2k+1)!} = x - \frac{x^3}{3!} + \frac{x^5}{5!} - \frac{x^7}{7!} + \cdots
\]

\begin{itemize}
    \item Escriba un programa que:
    \begin{itemize}
        \item Solicite al usuario un valor real \(x\).
        \item Solicite un entero positivo \(n\) (número de términos).
        \item Calcule la suma de los primeros \(n\) términos de la serie usando un ciclo \texttt{for}.
\item Compare el resultado con \texttt{math.sin(x)}.
\item Puede usar la función \texttt{factorial()} del módulo \texttt{math}.
    \end{itemize}
    
     \item Bonus: 
     \begin{itemize}
	\item Puede crear su propia función que implemente el factorial para obtener un bonus de décimas.
     \end{itemize}
\end{itemize}
\end{tcolorbox}

\begin{solution}

\subsubsection*{Idea}
La serie de Taylor de \(\sin(x)\) es una suma alternante de términos impares. Cada término tiene la forma:
\[
(-1)^k \frac{x^{2k+1}}{(2k+1)!}
\]

Para aproximar la función:
\begin{itemize}
    \item Se suman los primeros \(n\) términos de la serie.
    \item Se utiliza un ciclo \texttt{for} para iterar sobre los valores de \(k\).
    \item Se compara el resultado con el valor real usando \texttt{math.sin(x)}.
\end{itemize}

\subsubsection*{Implementación}

\begin{pythonjupyter}
import math

# Solicitar datos al usuario
x = float(input("Ingrese el valor de x: "))
n = int(input("Ingrese el número de términos: "))

# Inicializar suma
suma = 0

# Cálculo de la serie
for k in range(n):
    termino = ((-1)**k) * (x**(2*k + 1)) / math.factorial(2*k + 1)
    suma += termino

# Resultado aproximado
print(f"Aproximación de sin(x): {suma}")

# Valor real
valor_real = math.sin(x)
print(f"Valor real de sin(x): {valor_real}")

# Error absoluto
error = abs(valor_real - suma)
print(f"Error absoluto: {error}")
\end{pythonjupyter}

\subsubsection*{Resultado}
El programa aproxima \(\sin(x)\) mediante una serie de Taylor y permite comparar la precisión de la aproximación en función del número de términos utilizados.

\end{solution}
